%% toc: P2 23 23.1
\section{Church's simple theory of types: framing and further results}
\label{sec:p02-w02c-p02-028-church-supplement}

This section extends the account of Church's paper of 1940
\cite{Church1940} given in \cref{sec:p02-w02-p02-040-church}.  That section
fixes the type symbols, the well-formed formulas, the abbreviations, the
rules and the axioms of the system, and proves the selected theorems.  Here
we add four things the source states and that are not carried there: the
source's own reason for its formulation and its acknowledgements; the
source's examples of the type-symbol convention, with a worked
reconstruction; three qualifications that limit how the source's claims are
to be read; and proofs bearing on three claims of the source.

\emph{Inputs.}  The system of
\cref{def:p02-w02-p02-040-church-types,def:p02-w02-p02-040-church-wff,def:p02-w02-p02-040-church-abbrev,def:p02-w02-p02-040-church-rules,def:p02-w02-p02-040-church-axioms};
the theorems proved in \cref{sec:p02-w02-p02-040-church}; and, for the
semantic result, the standard models
\(\mathcal S(D,c)\) of \cref{def:p02-w02-p02-040-church-standard}, with ZFC
as metatheory.
\emph{Conclusions.}  (i) No formula free of the selection constants defines
the predecessor of the natural numbers of type \(\iota'\) within that type
(\cref{thm:p02-w02c-p02-028-predecessor}); this is the source's remark of
p.~59, proved from the interpretation with two individuals that refutes
Axiom~8.  (ii) The source's representation of classes by functions into the
numerals \(0_{\iota'},1_{\iota'}\) makes coextensive propositional functions
determine equal classes, using 10\(^{\iota'\alpha}\) and no extensionality
axiom for propositions (\cref{thm:p02-w02c-p02-028-classes}).  (iii) The
theorems 26\(^{\iota'}\), 26\(^{\iota''}\), \ldots{} follow without Axiom~7
once the extensionality axioms are present
(\cref{thm:p02-w02c-p02-028-no-ax7}); whether they follow from
Axioms 1--6\(^\alpha\) and 8 alone, as the source announces, is left open
here (\cref{def:p02-w02c-p02-028-open}).
\emph{Standing.}  Notation follows the source.  The precise formulations of
(i)--(iii) are the handbook's.  The source states (i) and the claim behind
(ii) without proof; for the claim behind (iii) its proof uses Axiom~7.

\subsection{The source's purpose and acknowledgements}

\begin{hbconvention}{def:p02-w02c-p02-028-purpose}{The purpose of the formulation}
The paper sets out to formulate the simple theory of types so that it takes
over certain features of the calculus of \(\lambda\)-conversion
\cite[p.~56]{Church1940}.  The incorporation is deliberately partial.  The
source holds that a complete incorporation is impossible as long as
\(\lambda x\) is to keep its meaning as an abstraction operator and
juxtaposition its meaning as application of a function to an argument
\cite[p.~56]{Church1940}.  The source offers the partial incorporation for
a positive reason of its own: it says that it ``has certain advantages from
the point of view of type theory'' \cite[p.~56]{Church1940}.  It also
detaches that reason from any verdict on type theory as a foundation: the
formulation is put forward as of interest on this basis whatever view one
takes of whether the theory of types is in the end a satisfactory
foundation for logic and mathematics \cite[p.~56]{Church1940}.  The paper
therefore argues for its formulation inside type theory and does not argue
for type theory itself.

In this handbook the ``features of \(\lambda\)-conversion'' are visible in
three places: terms of every type are formed by \(\lambda\)-abstraction
(\cref{def:p02-w02-p02-040-church-wff}); Rules I--III are the rules of
\(\lambda\)-conversion (\cref{def:p02-w02-p02-040-church-rules}); and
comprehension is obtained from abstraction and conversion instead of being
postulated (\cref{thm:p02-w02-p02-040-church-comprehension}).
\end{hbconvention}

\begin{hbconvention}{def:p02-w02c-p02-028-attributions}{The source's attributions}
The source records the following debts and historical claims; each is
reported here with the locator at which the source makes it.
\begin{enumerate}
 \item \emph{The simple theory of types.}  For the theory being formulated
   the source refers to Carnap's \emph{Abriss der Logistik} (Vienna 1929),
   \S9.  It adds that the simple theory was proposed as a modification of
   Russell's ramified theory of types by Leon Chwistek in 1921 and 1922 and by
   F.~P.~Ramsey in 1926 \cite[p.~56, n.~1]{Church1940}.
 \item \emph{The calculus of \(\lambda\)-conversion.}  For that calculus the
   source refers, as examples, to Church's mimeographed \emph{Mathematical
   logic} (Princeton, 1936) and to \emph{The calculi of lambda-conversion},
   then a forthcoming monograph \cite[p.~56, n.~2]{Church1940}.
 \item \emph{Features not tied to \(\lambda\)-conversion.}  For every feature
   of the formulation that is not directly connected with the incorporation
   of \(\lambda\)-conversion, the source acknowledges a heavy debt to
   Whitehead and Russell, to Hilbert and Ackermann, to Hilbert and Bernays,
   and to the forerunners of these works \cite[p.~56]{Church1940}.  Its
   references are: for Whitehead and Russell, Russell's ``Mathematical logic
   as based on the theory of types'', \emph{American Journal of Mathematics}
   30 (1908), pp.~222--262, and \emph{Principia Mathematica}, vol.~1 (1910;
   2nd ed.\ 1925), vol.~2 (1912; 2nd ed.\ 1927) and vol.~3 (1913; 2nd ed.\
   1927) \cite[p.~56, n.~3]{Church1940}; Hilbert and Ackermann,
   \emph{Grundz\"uge der theoretischen Logik} (Berlin 1928; 2nd ed.\ 1938)
   \cite[p.~56, n.~4]{Church1940}; and Hilbert and Bernays, \emph{Grundlagen
   der Mathematik}, vol.~1 (Berlin 1934) and vol.~2 (Berlin 1939)
   \cite[p.~56, n.~5]{Church1940}.
\end{enumerate}
These are the source's attributions.  The handbook has not checked the
cited works and makes no historical claim of its own about them.  The
source's division separates the features connected with
\(\lambda\)-conversion, among them abstraction and Rules I--III, which the
source calls the rules of \(\lambda\)-conversion
\cite[\S3, p.~60]{Church1940}, from the remaining features, for which it
acknowledges the works in item~3.  The source does not say which of the
remaining features comes from which work.
\end{hbconvention}

\subsection{The source's examples of the type-symbol convention}

\Cref{def:p02-w02-p02-040-church-types} states the convention that
parentheses are omitted from type symbols with association to the left
\cite[\S1, p.~56]{Church1940}.  The source illustrates it by three examples,
and the following reconstruction works through them.

\begin{hbexample}{def:p02-w02c-p02-028-type-examples}{The source's three type-symbol abbreviations}
The source gives \(o\iota\) for \((o\iota)\), \(\iota\iota\iota\) for
\(((\iota\iota)\iota)\), and \(\iota\iota(\iota\iota)\) for
\(((\iota\iota)(\iota\iota))\) \cite[\S1, p.~56]{Church1940}.  Restoring
parentheses means grouping the leftmost two components first and then
attaching each further component to the group so far.
\begin{enumerate}
 \item \(o\iota\) has two components, \(o\) and \(\iota\); one pair of
   parentheses gives \((o\iota)\), the type of functions from individuals to
   propositions, that is, of properties of individuals.
 \item \(\iota\iota\iota\) has three components.  Grouping the first two
   gives \((\iota\iota)\); attaching the third gives
   \(((\iota\iota)\iota)\).  By the reading of
   \cref{def:p02-w02-p02-040-church-types}, this is the type of functions
   whose arguments are individuals and whose values are of type
   \((\iota\iota)\), that is, curried two-place functions from individuals
   to individuals.
 \item \(\iota\iota(\iota\iota)\) has three components: \(\iota\),
   \(\iota\) and the parenthesized \((\iota\iota)\).  Grouping the first two
   gives \((\iota\iota)\); attaching the third gives
   \(((\iota\iota)(\iota\iota))\).  This is the type of functions from
   functions of type \((\iota\iota)\) to functions of type \((\iota\iota)\).
   It is also the type \(\iota'\) of the natural numbers of the source's
   arithmetic, since \(\alpha'\) abbreviates \(((\alpha\alpha)(\alpha\alpha))\)
   and here \(\alpha=\iota\).
\end{enumerate}
The parentheses inside \(\iota\iota(\iota\iota)\) cannot be dropped: by the
same rule \(\iota\iota\iota\iota\) would restore to
\((((\iota\iota)\iota)\iota)\), a different type.  The convention never
needs parentheses at the left end of a type symbol, only around a compound
argument type at the right.

\emph{Host form.}  Under the type map of
\cref{def:p02-w02-p02-040-church-bridge}, which reverses the order and
writes the argument type first, the three examples become
\(\iota\to o\), \(\iota\to(\iota\to\iota)\) and
\((\iota\to\iota)\to(\iota\to\iota)\); with \(\to\) associating to the right
the second is \(\iota\to\iota\to\iota\).  For instance
\(((\iota\iota)(\iota\iota))^\circ=(\iota\iota)^\circ\to(\iota\iota)^\circ
=(\iota\to\iota)\to(\iota\to\iota)\).
\end{hbexample}

The abbreviation \(\iota(\iota\iota)\) for \((\iota(\iota\iota))\) is also a
correct instance of the convention, but it is not one of the source's
examples, which are the three above; the source's third example is
\(\iota\iota(\iota\iota)\), with two components before the parenthesized
one.

\subsection{Three qualifications in the source}

\begin{hbwarning}{def:p02-w02c-p02-028-qualifications}{How the source limits its interpretive and syntactic claims}
\begin{enumerate}
 \item \emph{Interpretation is not part of the construction.}  After giving
   the intended reading of the types \(o\) and \(\iota\) and leaving their
   nature open, the source adds that interpretation is in any case irrelevant
   to the abstract construction of the theory, and that other, quite
   different interpretations are possible, with the parenthetical proviso
   ``formal consistency assumed'' \cite[\S1, p.~57]{Church1940}.  So the
   source makes no claim that the formal system is about propositions and
   individuals in any fixed sense.  The standard models of
   \cref{def:p02-w02-p02-040-church-standard} are the handbook's
   formalization of one intended interpretation, not part of the system;
   the consistency the proviso presupposes is proved in clause~2 of
   \cref{thm:p02-w02-p02-040-church-soundness}.
 \item \emph{Why formal and syntactical theorems are distinguished.}  The
   source distinguishes theorems of the system from theorems about the
   system; the deduction theorem
   (\cref{thm:p02-w02-p02-040-church-deduction}) is of the second kind.
   Its stated reason is that
   such distinctions are a necessary part of using one language that is
   already known, English, to set up another, more exact language
   \cite[\S4, p.~61]{Church1940}.  The distinction is therefore presented as
   a consequence of the method of construction, not as a feature of the
   particular system.
 \item \emph{The status of 21\(^{\alpha\beta}\).}  The source lists
   21\(^{\alpha\beta}\), \(f_{\alpha\beta}=\lambda x_\beta(f_{\alpha\beta}x_\beta)\),
   as a consequence of Axioms 1--4 and 10\(^{\alpha\beta}\), and says at once
   that no use will be made of these theorems below, because the rest of the
   paper is concerned only with consequences of Axioms 1--9\(^\alpha\)
   \cite[\S5, p.~63]{Church1940}.  So none of the source's later derivations
   rests on the \(\eta\)-equation.  The handbook proves 21\(^{\alpha\beta}\)
   in \cref{thm:p02-w02-p02-040-church-12-21}; of the results below,
   \cref{thm:p02-w02c-p02-028-classes,thm:p02-w02c-p02-028-no-ax7}
   use Axiom~10 directly and do not use 21.
\end{enumerate}
\end{hbwarning}

\subsection{The same-type predecessor and descriptions}

The source defines the predecessor \(P_{\alpha'\alpha'}\) within one type
through \(T_{\alpha''\alpha'}\), and \(T_{\alpha''\alpha'}\) is a description
(\cref{def:p02-w02-p02-040-church-abbrev}).  Kleene's formula
\(P_{\alpha'\alpha'''}\), by contrast, uses no selection constant, but it
lowers the type: from a number of type \(\alpha'''\) it yields a number of
type \(\alpha'\).  The source says that the same-type predecessor cannot be
defined without descriptions and that this follows from the independence of
the axiom of infinity \cite[\S2, p.~59]{Church1940}.  The following theorem
gives the claim a precise form and proves it by the interpretation with two
individuals, which is the interpretation the source uses to show Axiom~8
independent \cite[\S4, p.~61]{Church1940}.

\begin{hbtheorem}{thm:p02-w02c-p02-028-predecessor}{proved=annex:p02-w02c-p02-028-predecessor}{No description-free same-type predecessor}
Let \(\mathbf F_{\iota'\iota'}\) be a wff with no free variables in which no
selection constant \(\iota_{\beta(o\beta)}\) occurs, for any type
\(\beta\).  Let \(D\) be a set with at least two members and \(c\) a family
of choice functions as in \cref{def:p02-w02-p02-040-church-standard}.  Then
in \(\mathcal S(D,c)\)
\[
 \sem{\mathbf F_{\iota'\iota'}(S_{\iota'\iota'}0_{\iota'})}\ne\sem{0_{\iota'}}
 \quad\text{or}\quad
 \sem{\mathbf F_{\iota'\iota'}(S_{\iota'\iota'}2_{\iota'})}\ne\sem{2_{\iota'}}.
\]
Consequently the formula
\[
 N_{o\iota'}x_{\iota'}\supset\mathbf F_{\iota'\iota'}(S_{\iota'\iota'}x_{\iota'})=x_{\iota'}
\]
is not valid in \(\mathcal S(D,c)\), and it is not a theorem of the system
with all of Axioms 1--11.
\end{hbtheorem}

The formula \(\mathbf F_{\iota'\iota'}\) may contain \(N_{oo}\),
\(A_{ooo}\) and the quantifier constants \(\Pi_{o(o\beta)}\) at any types,
and abstraction of any form.  A definition of the same-type predecessor
\(P_{\alpha'\alpha'}\) by typical ambiguity would include the case
\(\alpha=\iota\); so no single description-free schema defines it at every
type.  The proof shows where descriptions matter: it compares
\(\mathcal S(D,c)\) with the model on two individuals through a relation
between the two models that every description-free wff respects, and the
selection constants are the only primitive symbols for which that step can
fail.  The theorem does not show that the source's own
\(P_{\alpha'\alpha'}\) denotes the predecessor; that is the source's stated
intention \cite[\S2, p.~59]{Church1940}, which the handbook reports and does
not prove.

\subsection{Classes as functions into the numerals}

The source argues that Axioms 10\(^{\alpha\beta}\), though weaker in some
directions than other extensionality axioms, are adequate, because a class
can be represented by a function into the numerals of type \(\iota'\)
\cite[\S4, p.~61]{Church1940}; the representing formula is displayed in
\cref{sec:p02-w02-p02-040-church}.  The point of the representation is
that \(10^{\alpha\beta}\) gives extensionality for functions but not for
propositions, so that two coextensive propositional functions need not be
provably equal; the source notes that the further axiom
\(p\equiv q\supset p=q\) would close that gap at the price of leaving only
two propositions \cite[\S4, p.~61]{Church1940}.  The theorem below proves
that the representing functions are extensional without that further axiom.

\begin{hbdefinition}{def:p02-w02c-p02-028-class}{The class associated with a propositional function}
Let \(\mathbf F_{o\alpha}\) be a wff, and let \(x_\alpha\) and
\(y_{\iota'}\) be variables not free in it.  The \emph{class associated with}
\(\mathbf F_{o\alpha}\) is the wff
\[
 \mathbf C_{\mathbf F}\ :=\ \lambda x_\alpha(\text{\rotatebox[origin=c]{180}{\(\iota\)}}y_{\iota'})\,.\,
 (\mathbf F_{o\alpha}x_\alpha)[y_{\iota'}=0_{\iota'}]\vee
 (\sim\mathbf F_{o\alpha}x_\alpha)[y_{\iota'}=1_{\iota'}]
\]
of type \(\iota'\alpha\) \cite[\S4, p.~61]{Church1940}, in which the
juxtapositions inside the description are conjunctions.  The name
\(\mathbf C_{\mathbf F}\) is the handbook's.
\end{hbdefinition}

\begin{hbtheorem}{thm:p02-w02c-p02-028-classes}{proved=annex:p02-w02c-p02-028-classes}{Coextensive propositional functions determine equal classes}
Let \(\mathbf F_{o\alpha}\) and \(\mathbf G_{o\alpha}\) be wffs in which
\(x_\alpha\) and \(y_{\iota'}\) are not free.
\begin{enumerate}
 \item Using Axioms 1--6\(^\gamma\), 9\(^{\iota'}\) and
   10\(^{\iota'\alpha}\) only,
   \(\derives(x_\alpha)[\mathbf F_{o\alpha}x_\alpha\equiv\mathbf G_{o\alpha}x_\alpha]
   \supset\mathbf C_{\mathbf F}=\mathbf C_{\mathbf G}\).
 \item Using Axioms 1--6\(^\gamma\), 7 and 9\(^{\iota'}\) only,
   \(\derives\mathbf C_{\mathbf F}x_\alpha=0_{\iota'}\equiv\mathbf F_{o\alpha}x_\alpha\)
   and
   \(\derives\mathbf C_{\mathbf F}=\mathbf C_{\mathbf G}\supset
   (x_\alpha)[\mathbf F_{o\alpha}x_\alpha\equiv\mathbf G_{o\alpha}x_\alpha]\).
\end{enumerate}
\end{hbtheorem}

Clause~1 is the adequacy the source claims: equality of classes, which is
equality of functions of type \(\iota'\alpha\), follows from coextensiveness
by an axiom of extensionality for functions alone.  Clause~2 shows that
membership, read as \(\mathbf C_{\mathbf F}x_\alpha=0_{\iota'}\), recovers
\(\mathbf F_{o\alpha}\), and that equal classes have coextensive
propositional functions.  Clause~2 needs \(0_{\iota'}\ne1_{\iota'}\), which
comes from 28\(^\iota\) and so from Axiom~7.  It cannot be dropped: in the
standard model with one individual, the type \(\iota'\) has a single member,
so \(0_{\iota'}\) and \(1_{\iota'}\) have the same denotation and so do all
the \(\mathbf C_{\mathbf F}\) for a given \(\alpha\); the first formula of
clause~2 is then false whenever \(\mathbf F_{o\alpha}x_\alpha\) is false,
while Axioms 1--6\(^\gamma\) and 9\(^{\iota'}\) are valid there
(\cref{thm:p02-w02-p02-040-church-soundness}).

\subsection{Theorem 26 at the types \(\iota',\iota'',\ldots\) without Axiom~7}

The source announces that 26\(^{\iota'}\), 26\(^{\iota''}\), \ldots{} are
consequences of Axioms 1--6\(^\alpha\) and 8 \cite[\S6, p.~64]{Church1940}.
The derivation of \cref{thm:p02-w02-p02-040-church-26-30}, which follows the
source's proof, also uses Axiom~7, through 28\(^\alpha\) and 27\(^\iota\)
(\cref{def:p02-w02-p02-040-church-26-axiom-7}).  The next theorem removes
Axiom~7 at the cost of extensionality axioms.  Put \(\alpha_0=\iota\) and
\(\alpha_{k+1}=\alpha_k'\).

\begin{hbtheorem}{thm:p02-w02c-p02-028-no-ax7}{proved=annex:p02-w02c-p02-028-no-ax7}{Theorem 26 at higher types without Axiom~7}
For every \(k\ge1\), the formula 26\(^{\alpha_k}\) is derivable from Axioms
1--6\(^\gamma\), Axiom~8, and the instances 10\(^{\alpha_j\alpha_j}\) and
10\(^{(\alpha_j\alpha_j)(\alpha_j\alpha_j)}\) for \(j\le k\), without
Axiom~7.
\end{hbtheorem}

\begin{hbwarning}{def:p02-w02c-p02-028-open}{What remains open about the source's announcement}
\Cref{thm:p02-w02c-p02-028-no-ax7} uses extensionality, and the
source's announcement does not mention it.  Whether 26\(^{\iota'}\),
26\(^{\iota''}\), \ldots{} are derivable from Axioms 1--6\(^\alpha\) and 8
alone, with neither Axiom~7 nor Axiom~10, is not settled here.  The proof of
\cref{thm:p02-w02c-p02-028-no-ax7} splits on 27\(^\iota\).  When
27\(^\iota\) holds, the source's route applies.  When it fails, Axiom~10
shows that every type \(\alpha_k\) has at most one member, so 26 holds
trivially.  Without Axiom~10, the case with one individual gives no
information about the higher types, and the source's route is blocked:
it needs 28\(^\iota\), and 28\(^\iota\) is false in the standard model with
one individual, where Axiom~8 is valid.  So 28\(^\iota\) is not derivable
from Axioms 1--6\(^\alpha\) and 8, and a proof of the announcement without
Axiom~10 would have to avoid it.
\end{hbwarning}

\subsection{Relationships and sources}

The results here extend \cref{sec:p02-w02-p02-040-church}.
\Cref{thm:p02-w02c-p02-028-predecessor} gives the remark
reported in \cref{def:p02-w02-p02-041-church-predecessor} a precise form
and a proof, and it uses the independence argument of clause~3 of
\cref{thm:p02-w02-p02-040-church-soundness}.
\Cref{thm:p02-w02c-p02-028-classes} supplies the derivation behind the
class representation discussed after \cref{thm:p02-w02-p02-040-church-comprehension}.
\Cref{thm:p02-w02c-p02-028-no-ax7} bears on
\cref{def:p02-w02-p02-040-church-26-axiom-7} and on the footnote-12
questions recorded in \cref{def:p02-w02b-p02-034-church-open}, none of
which it settles.  The source throughout is \cite{Church1940}, pp.~56--68.
